\section{Auditoría de fuentes y pregunta central del curso}

Esta clase no anuncia que «la robótica ya tuvo su momento GPT», sino que propone una receta de investigación de 2023: ¿es posible combinar los principios de diseño del scaling en aprendizaje automático, los Internet-scale language / vision models en constante avance y los datos robóticos fuera de línea a gran escala, de modo que las capacidades de los robots reales dejen de crecer lentamente a fuerza de ingeniería manual, una por una? A continuación, Ted Xiao usa RT-1, SayCan, Inner Monologue y DIAL para mostrar cómo esta receta entra, respectivamente, en skill learning, planning, feedback y data augmentation.

\subsection{Cómo se reconstruyó esta clase}

El sitio oficial del curso ofrece la grabación y tres artículos de lectura obligatoria, pero no publica un slide deck independiente. Por ello, esta clase toma como base la grabación oficial de Stanford Online en 1080p: se muestreó la región de screen-share cada dos segundos, de 2,284 cuadros se seleccionaron 315 candidatos estables de alta exhaustividad y, tras eliminar duplicados manualmente, quedaron 45 estados didácticos independientes; por último, se volvieron a extraer cuadros de ancho completo del video original de 1920x1080 según las marcas de tiempo verificadas, para evitar que se recortara el lado derecho de las diapositivas. La motivación expresada oralmente y las preguntas y respuestas provienen de los subtítulos humanos oficiales `en-US`; los subtítulos locales anteriores eran 4,183 subtítulos automáticos con desplazamiento repetido, y se sustituyeron por 1,840 captions humanos no vacíos.

\begin{longtable}{L{0.21\textwidth}L{0.31\textwidth}L{0.38\textwidth}}
\toprule
Capa de evidencia & Material local & Función en la nota \\
\midrule
Grabación oficial & Video de 1:16:07 y 45 teaching states recuperados & Determina el orden de la clase, los diagramas de arquitectura, las tablas de resultados, las demostraciones y el mapa del sistema \\
Subtítulos humanos oficiales & `lecture15.en.srt` y timed transcript & Recupera la speaker motivation, las decisiones de ingeniería, las limitaciones señaladas en las preguntas y la incertidumbre \\
Artículos obligatorios del curso & RT-1, SayCan, Inner Monologue & Verifica fórmulas, definiciones experimentales, alcance de los resultados e interfaces entre módulos \\
Artículos primarios complementarios & BC-Z y DIAL & Explican los antecedentes históricos, los datos fuera de línea y el VLM relabeling \\
Archivos de auditoría & manifest, coverage, teacher-voice ledger & Demuestran que cada elemento visual y oral tiene un destino didáctico \\
\bottomrule
\end{longtable}

\begin{warningbox}{Límite histórico}
La clase se impartió el 7 de febrero de 2023. RT-2, Open X-Embodiment, PaLM-E, Gemini Robotics y los resultados posteriores de vision-language-action scaling no pueden proyectarse hacia atrás como evidencia de la clase. Al final se incluye por separado una interpretación de los VLA modernos, pero se mantiene separada de lo que Ted Xiao sostenía en ese momento.
\end{warningbox}

\subsection{Primero, distinguir términos cercanos}

El ponente usa oralmente foundation model, Internet-scale model y LLM de manera intercambiable y poco estricta, pero al escribir la nota es necesario distinguir el objeto de entrenamiento, la fuente de datos y la función dentro del sistema. El siguiente glosario establece primero las coordenadas de lectura; las fórmulas y los casos vienen más adelante.

\begin{longtable}{L{0.22\textwidth}L{0.30\textwidth}L{0.38\textwidth}}
\toprule
Término & Problema que resuelve & Mecanismo central y relación con el curso \\
\midrule
Foundation model & Si un solo modelo puede reutilizarse en una gran cantidad de tareas posteriores & Depende del entrenamiento a escala, una interfaz general y representaciones transferibles; esta clase lo plantea como objetivo y no como un hecho consumado \\
Imitation learning & Cómo aprender una policy a partir de demostraciones cuando no hay una reward explícita & Aprendizaje supervisado $p(a\mid o,l)$; el punto de partida estable de RT-1 \\
Offline robot learning & Cómo desacoplar la recolección del entrenamiento & Primero se recolecta un conjunto de datos fijo y luego se entrenan varios métodos sobre él repetidamente; reduce el ensayo y error en línea con robots reales \\
Tokenización & Cómo hacer que imágenes, lenguaje y acciones entren en un modelo de secuencias unificado & Comprime entradas continuas o de alta dimensión en tokens discretos o compactos; la interfaz de interoperabilidad de RT-1 \\
Affordance & Si una habilidad es ejecutable en el estado actual & SayCan usa un value / success score para acotar el conocimiento previo lingüístico del LLM \\
Closed-loop planning & Cómo actualizar el plan cuando falla la ejecución & Inner Monologue retroalimenta al planner con scene, clarification y success feedback \\
Vision-language model & Cómo reescribir robot data con Internet-scale visual semantics & DIAL usa un VLM para generar etiquetas de lenguaje más ricas, en lugar de controlar directamente al robot \\
\bottomrule
\end{longtable}

\subsection{Cuatro preguntas que atraviesan toda la clase}

Al leer los sistemas siguientes conviene preguntarse continuamente: si a los datos robóticos les falta cantidad o diversidad; si el valor del Transformer proviene del número de parámetros o de una token interface unificada; cómo se acotan las acciones propuestas por el LLM según su ejecutabilidad física; y cómo la retroalimentación y las etiquetas de lenguaje modifican la siguiente ronda de la policy. Solo al ubicar las cuatro preguntas en una misma cadena del sistema se ve con claridad que los cuatro trabajos no son demos inconexas.

\begin{importantbox}{La tesis sistémica de esta clase}
La dificultad de un foundation model para robótica no consiste solo en entrenar una red más grande, sino en dotar de interfaces combinables a los datos de habilidades, el control en tiempo real, la planificación en lenguaje, la retroalimentación del entorno y las etiquetas semánticas. En esta recipe de 2023, el lenguaje funciona más bien como «universal glue»: conecta los módulos, pero no puede sustituir el grounding físico de cada uno.
\end{importantbox}

\subsection{Resumen de la sección}

Esta clase se reconstruyó con la grabación oficial, los subtítulos humanos y cinco artículos primarios, y distingue estrictamente entre una «recipe de investigación» y un «foundation model ya demostrado». A continuación se explica primero por qué la robótica aspira a la emergence y la homogenization, y después por qué el mundo físico se ha convertido en la última ficha de dominó difícil de derribar.
